% PhD Defense Slides — Daniel Abreu
% Three Essays on Factor Models
% ISEG – Lisbon School of Economics and Management, Universidade de Lisboa
% ~28 slides, 30-minute presentation

\documentclass{beamer}

% Packages
\usepackage{amsmath,amssymb,bm}
\usepackage{booktabs}
\usepackage{graphicx}
\usepackage{natbib}
\usepackage{lmodern}

% Style: match existing Oxford presentation
\usefonttheme{professionalfonts}
\usetheme{default}
\setbeamertemplate{navigation symbols}{}
\setbeamertemplate{footline}{%
  \leavevmode\hbox{%
    \begin{beamercolorbox}[wd=.5\paperwidth,ht=2.25ex,dp=1ex,left]{author in head/foot}%
      \hspace{1em}\insertshortauthor
    \end{beamercolorbox}%
    \begin{beamercolorbox}[wd=.5\paperwidth,ht=2.25ex,dp=1ex,right]{title in head/foot}%
      \insertframenumber{} / \inserttotalframenumber\hspace{1em}
    \end{beamercolorbox}}%
}

\definecolor{isegblue}{RGB}{40,60,130}
\setbeamercolor{frametitle}{fg=isegblue}
\setbeamerfont{frametitle}{size=\large,series=\bfseries}
\setbeamercolor{structure}{fg=isegblue}

% Figure path
\graphicspath{%
  {../Thesis_20251001/Smooth_transition_factor_model_paper/Chapters/}%
  {../Thesis_20251001/Non-stationary dynamic factor model with threshold cointegration/}%
  {../Thesis_20251001/Financial cycle/}%
}

% Commands
\newcommand{\F}{\bm{F}}
\newcommand{\X}{\bm{X}}
\newcommand{\La}{\bm{\Lambda}}
\newcommand{\eps}{\bm{\varepsilon}}

% Title
\title[\textcolor{isegblue}{Three Essays on Factor Models}]{Three Essays on Factor Models}
\author[Daniel Abreu]{Daniel Abreu}
\institute{ISEG -- Lisbon School of Economics and Management\\
Universidade de Lisboa\\[0.5em]
{\small Supervisors: Prof.\ João Nicolau \& Prof.\ Paulo Rodrigues}}
\date{PhD Defense, 2025}

% =========================================================
\begin{document}
% =========================================================

% ----------------------------------------------------------
% SLIDE 1: Title
% ----------------------------------------------------------
\begin{frame}
  \titlepage
\end{frame}

% ----------------------------------------------------------
% SLIDE 2: Thesis Overview
% ----------------------------------------------------------
\begin{frame}{Overview: Three Essays on Factor Models}

\begin{block}{Common Thread}
  High-dimensional datasets are ubiquitous in economics and finance.\\
  This thesis extends factor models to accommodate \textbf{nonlinearity} and \textbf{nonstationarity}.
\end{block}

\vspace{0.6em}

{\small
\begin{tabular}{l p{4.2cm} p{3.8cm}}
\toprule
 & \textbf{Essay} & \textbf{Key contribution} \\
\midrule
\textcolor{isegblue}{\textbf{I}} & Smooth transition factor model & Gradual regime shifts in loadings \\[4pt]
\textcolor{isegblue}{\textbf{II}} & Nonstationary DFM + threshold & Nonlinear cointegration in factors \\[4pt]
\textcolor{isegblue}{\textbf{III}} & Financial cycle measurement & DFM-based early-warning indicator \\
\bottomrule
\end{tabular}
}

\vspace{0.6em}

{\small All three essays contribute to the growing literature on factor models for large panels.}

\end{frame}

% ==========================================================
% PART I — SMOOTH TRANSITION FACTOR MODEL
% ==========================================================

\begin{frame}{}
  \begin{center}
    \Large\textcolor{isegblue}{\textbf{Essay I}} \\[0.8em]
    \large Smooth Transition Factor Model
  \end{center}
\end{frame}

% ----------------------------------------------------------
% SLIDE 3: STFM Motivation
% ----------------------------------------------------------
\begin{frame}{Essay I: Motivation}

\textbf{Standard factor model:}
\[ \X_t = \La \F_t + \eps_t \]
Factor loadings $\La$ are \emph{assumed constant} across time.

\vspace{0.8em}
\textbf{But economic relationships change over time:}
\begin{itemize}
  \item Loading of credit variables shifts across business cycle phases
  \item Abrupt structural-break models impose an all-or-nothing regime switch
  \item In practice, transitions are \textbf{gradual} and \textbf{continuous}
\end{itemize}

\vspace{0.8em}
\textbf{This paper:} Allow loadings to vary \emph{smoothly} between two regimes using a logistic transition function.

\end{frame}

% ----------------------------------------------------------
% SLIDE 4: STFM Model Setup
% ----------------------------------------------------------
\begin{frame}{Essay I: Model}

\textbf{Smooth Transition Factor Model (STFM):}
\[
x_{it} = \bm{\lambda}_{1i}' \F_t \,+\, \bm{\lambda}_{2i}' \F_t \cdot G(s_t;\,\gamma,c) \,+\, \varepsilon_{it}
\]

\begin{itemize}
  \item $\F_t$: $r \times 1$ vector of common factors
  \item $\bm{\lambda}_{1i}, \bm{\lambda}_{2i}$: regime-specific loadings
  \item $G(s_t;\gamma,c) = \bigl[1 + e^{-\gamma(s_t - c)}\bigr]^{-1} \in (0,1)$: logistic transition
  \item $s_t$: transition variable (e.g., lagged output growth); $\gamma > 0$: speed; $c$: location
\end{itemize}

\vspace{0.6em}
\textbf{Interpretation:}
\begin{itemize}
  \item When $G \approx 0$: loadings $\approx \bm{\lambda}_{1i}$ (regime 1, e.g., expansion)
  \item When $G \approx 1$: loadings $\approx \bm{\lambda}_{1i} + \bm{\lambda}_{2i}$ (regime 2, e.g., recession)
  \item Transition is \textbf{smooth}: both regimes can coexist at any point in time
\end{itemize}

\end{frame}

% ----------------------------------------------------------
% SLIDE 5: STFM Estimation
% ----------------------------------------------------------
\begin{frame}{Essay I: Estimation}

\textbf{Two-step profiled NLS estimator:}

\begin{enumerate}
  \item \textbf{Profile over $(\gamma, c)$:} For each candidate $(\gamma, c)$ on a grid,\\
    construct $\widetilde{\X}_t(\gamma,c) = [X_t,\, X_t \cdot G(s_t;\gamma,c)]$ and extract factors via PCA.
  \item \textbf{Select $(\hat{\gamma}, \hat{c})$:} Minimize the sum of squared residuals across the grid.
  \item \textbf{Factor number} $r$: information criterion adapted for regime-shift models.
\end{enumerate}

\vspace{0.8em}
\textbf{Key theoretical results:}
\begin{itemize}
  \item \textbf{Consistency} of $\hat{\gamma}$, $\hat{c}$ as $N,T \to \infty$
  \item \textbf{Linearity test}: LM statistic based on Taylor-series expansion of $G$ around $\gamma = 0$
\end{itemize}

\end{frame}

% ----------------------------------------------------------
% SLIDE 6: STFM Linearity Test
% ----------------------------------------------------------
\begin{frame}{Essay I: Linearity Test}

\textbf{Problem:} Testing $H_0: \gamma = 0$ (standard DFM) is non-standard — $c$ is a nuisance parameter unidentified under $H_0$.

\vspace{0.6em}
\textbf{Solution:} Taylor-expand $G(s_t;\gamma,c)$ around $\gamma = 0$:
\[
G(s_t;\gamma,c) \approx \tfrac{1}{2} + \tfrac{\gamma}{4}(s_t - c) + O(\gamma^3)
\]
Under $H_0$: loading becomes $\bm{\lambda}_{1i} + \tfrac{1}{2}\bm{\lambda}_{2i}$ (constant) — the nuisance parameter $c$ drops out.

\vspace{0.6em}
\textbf{LM statistic:} Test $H_0: \bm{\delta}_i = 0$ in the auxiliary regression
\[
\hat{\varepsilon}_{it} = \bm{\lambda}_{1i}' \F_t + \bm{\delta}_{i}' (\F_t \cdot s_t) + u_{it}
\]
Asymptotically $\chi^2$ distributed under $H_0$.

\end{frame}

% ----------------------------------------------------------
% SLIDE 7: STFM Monte Carlo
% ----------------------------------------------------------
\begin{frame}{Essay I: Monte Carlo Evidence}

\begin{itemize}
  \item DGPs: STFM with $r = 1, 2, 3$ factors; $N \in \{50, 100\}$; $T \in \{100, 200\}$
  \item Transition variable $s_t$: NBER recession indicator (business cycle timing)
\end{itemize}

\vspace{0.4em}
\textbf{Main findings:}
\begin{itemize}
  \item \textbf{Bias and RMSE} for $\hat{\gamma}$ and $\hat{c}$ decrease as $N, T$ grow — consistent
  \item \textbf{Factor estimates} track DGP factors closely; space recovery improves with $N$
  \item \textbf{Linearity test}: correct size under $H_0$; power increases sharply with $T$
  \item \textbf{IC factor selection}: selects correct $r$ with high probability in moderate samples
\end{itemize}

\vspace{0.4em}
\textit{Estimation performs well for moderate-sized panels typical in macroeconomics.}

\end{frame}

% ----------------------------------------------------------
% SLIDE 8: STFM Empirical Application
% ----------------------------------------------------------
\begin{frame}{Essay I: Empirical Application — FRED-QD}

\textbf{Data:} 127 US macroeconomic quarterly series from FRED-QD.\\
\textbf{Transition variable:} lagged real GDP growth.

\vspace{0.5em}
\textbf{Model selection:} $\hat{r} = 3$ factors; linearity strongly rejected ($p < 0.01$).

\vspace{0.5em}
\textbf{Transition over time:}
\begin{center}
  \includegraphics[width=0.7\textwidth]{stfm_transition_over_time.pdf}
\end{center}

{\small Transition function tracks NBER recessions (shaded); factors shift smoothly during downturns.}

\end{frame}

% ----------------------------------------------------------
% SLIDE 9: STFM Forecasting & Takeaway
% ----------------------------------------------------------
\begin{frame}{Essay I: Forecasting Results \& Takeaway}

\textbf{Out-of-sample forecasting} (1-step ahead, rolling window, 1990--2019):

\vspace{0.4em}
\begin{tabular}{lcc}
\toprule
 & \textbf{STFM} & Linear DFM \\
\midrule
Industrial production & \textbf{0.87} & 1.00 \\
Manufacturers' orders & \textbf{0.81} & 1.00 \\
\bottomrule
\end{tabular}

{\small RMSE relative to linear DFM; below 1 = improvement.}

\vspace{0.8em}
\textbf{Takeaway:}
\begin{itemize}
  \item Smooth regime shifts in loadings are empirically relevant in US macro data
  \item The STFM provides a tractable, consistent framework for modeling gradual transitions
  \item Useful for \textbf{nowcasting} and \textbf{current economic assessment}
\end{itemize}

\end{frame}

% ==========================================================
% PART II — NON-STATIONARY DFM WITH THRESHOLD COINTEGRATION
% ==========================================================

\begin{frame}{}
  \begin{center}
    \Large\textcolor{isegblue}{\textbf{Essay II}} \\[0.8em]
    \large Nonstationary Dynamic Factor Model\\
    with Threshold Cointegration
  \end{center}
\end{frame}

% ----------------------------------------------------------
% SLIDE 10: NSTFM Motivation
% ----------------------------------------------------------
\begin{frame}{Essay II: Motivation}

\textbf{Standard DFM assumes stationarity} — but many macroeconomic and financial series are I(1):
\begin{itemize}
  \item Interest rates, exchange rates, asset prices, credit aggregates
\end{itemize}

\vspace{0.6em}
\textbf{Two existing approaches have limitations:}
\begin{itemize}
  \item \textit{First-difference everything}: discards long-run information
  \item \textit{Stock \& Watson (2016)}: allows I(1) factors but assumes \textbf{linear} adjustment
\end{itemize}

\vspace{0.6em}
\textbf{This paper:} Combine nonstationary two-level factor models with \textbf{threshold cointegration}.

\vspace{0.6em}
\textbf{Motivation from theory:}
\begin{itemize}
  \item Adjustment costs and financial frictions imply \textbf{inaction bands}
  \item Error-correction should be \textbf{inactive} under small deviations, \textbf{active} under large ones
\end{itemize}

\end{frame}

% ----------------------------------------------------------
% SLIDE 11: NSTFM Model Setup
% ----------------------------------------------------------
\begin{frame}{Essay II: Two-Level Factor Structure}

\textbf{Two-level nonstationary factor model:}

Group-level observation equation:
\[
x_{it} = \bm{\lambda}_{il}' \F_{lt} + \varepsilon_{it}, \qquad i \in \text{group } l
\]

Global-level linkage (both I(1)):
\[
\F_{lt} = \bm{A}_l \, G_t + \bm{\eta}_{lt}
\]

\begin{itemize}
  \item $G_t$: $r_G \times 1$ I(1) global factors
  \item $\F_{lt}$: $r_l \times 1$ I(1) group-specific factors
  \item Both levels may be \textbf{cointegrated}
\end{itemize}

\vspace{0.5em}
\textbf{Key insight:} The factors $G_t$ and $\F_{lt}$ share common stochastic trends. Modeling this linkage correctly is essential for consistent estimation.

\end{frame}

% ----------------------------------------------------------
% SLIDE 12: NSTFM Threshold & Band VECM
% ----------------------------------------------------------
\begin{frame}{Essay II: Threshold and Band VECM}

\textbf{Error-correction mechanism between factor levels:}

\vspace{0.4em}
\textbf{Threshold factor VECM} (asymmetric adjustment):
\[
\Delta \F_t = \alpha_1 \bm{\beta}' \F_{t-1} \cdot \mathbf{1}(q_t > 0) \;+\; \alpha_2 \bm{\beta}' \F_{t-1} \cdot \mathbf{1}(q_t \leq 0) \;+\; u_t
\]

\textbf{Band factor VECM} (no adjustment within band):
\[
\Delta \F_t = \alpha \bm{\beta}' \F_{t-1} \cdot \mathbf{1}(|\bm{\beta}' \F_{t-1}| > d) \;+\; u_t
\]

\begin{itemize}
  \item $\bm{\beta}$: cointegrating vector; $q_t = \bm{\beta}' \F_{t-1}$: equilibrium error
  \item $d > 0$: threshold parameter (band width)
  \item \textbf{Band model}: error correction is \textbf{inactive} when deviations are small (inside band)
\end{itemize}

\end{frame}

% ----------------------------------------------------------
% SLIDE 13: NSTFM Estimation
% ----------------------------------------------------------
\begin{frame}{Essay II: Estimation Procedure}

\textbf{Three-stage procedure:}

\begin{enumerate}
  \item \textbf{Initialization}: SVD of the stacked panel $\bm{X}$ to obtain starting values for $\F$ and $\bm{\Lambda}$
  \item \textbf{Joint grid search}: For each candidate threshold $d$ and cointegrating vector $\bm{\beta}$:
    \begin{itemize}
      \item Estimate VECM adjustment coefficients via OLS/RR
      \item Retain $(\hat{\bm{\beta}}, \hat{d})$ minimizing the objective function
    \end{itemize}
  \item \textbf{Sequential least squares}: Update $\F$, $\bm{\Lambda}$, $\bm{A}$ iteratively until convergence
\end{enumerate}

\vspace{0.5em}
\textbf{Reduced-rank regression} applied at the cointegrating step to impose the cointegration rank constraint.

\end{frame}

% ----------------------------------------------------------
% SLIDE 14: NSTFM Monte Carlo
% ----------------------------------------------------------
\begin{frame}{Essay II: Monte Carlo Evidence}

\textbf{Design:} Band factor VECM DGP with known $(\bm{\beta}, d)$.\\
$N \in \{20, 50, 100\}$, $T \in \{100, 200, 400\}$, 500 replications.

\vspace{0.4em}
\textbf{Key results:}
\begin{itemize}
  \item \textbf{Threshold estimator} $\hat{d}$: low bias, concentrated distributions; improves with $T$
  \item \textbf{Factor space recovery}: trace statistic decreases toward 0 as $N,T \to \infty$
  \item \textbf{Cointegrating vector}: estimates cluster tightly around true $\bm{\beta}$ even for $T = 100$
  \item Performance is \textbf{adequate} for datasets of the size typical in international macro (e.g., $N = 20$--$50$ countries, $T = 150$--$300$ months)
\end{itemize}

\end{frame}

% ----------------------------------------------------------
% SLIDE 15: NSTFM Empirical Application
% ----------------------------------------------------------
\begin{frame}{Essay II: Global Term Structure of Interest Rates}

\textbf{Data:} Government bond yields for 22 OECD countries, 3 maturities (2, 5, 10Y), 1985--2022.\\
\textbf{Structure:} Two-level — global yield factor + country-group factors.

\vspace{0.5em}
\textbf{Band VECM results:}
\begin{itemize}
  \item \textbf{Cointegration confirmed}: global and group factors share a common trend
  \item \textbf{Threshold $\hat{d}$}: estimated band of approximately ±60 bps
  \item \textbf{Inside band}: error-correction \emph{inactive} — deviations persist without mean reversion
  \item \textbf{Outside band}: error-correction \emph{activates} strongly ($\hat{\alpha} < 0$, significant)
\end{itemize}

\vspace{0.4em}
\textbf{Interpretation:} Financial market frictions and policy inertia create \textbf{inaction zones} — yield differentials are tolerated up to a threshold, but large deviations trigger active rebalancing.

\end{frame}

% ----------------------------------------------------------
% SLIDE 16: NSTFM Takeaway
% ----------------------------------------------------------
\begin{frame}{Essay II: Key Takeaway}

\begin{itemize}
  \item Standard linear cointegration \textbf{misspecifies} adjustment in financial markets
  \item The \textbf{band factor VECM} provides a rigorous framework for nonlinear long-run adjustment in large panels of nonstationary series
  \item \textbf{Estimation}: tractable, consistent, performs well in moderate samples
  \item \textbf{Empirical finding}: global interest rate co-movements are governed by threshold dynamics consistent with \textbf{transaction costs} and \textbf{policy inertia} models
  \item Opens the door to \textbf{impulse response analysis} and \textbf{regime-conditional forecasting} under nonstationary factors
\end{itemize}

\end{frame}

% ==========================================================
% PART III — FINANCIAL CYCLE MEASUREMENT
% ==========================================================

\begin{frame}{}
  \begin{center}
    \Large\textcolor{isegblue}{\textbf{Essay III}} \\[0.8em]
    \large Measuring the Financial Cycle:\\
    A Dynamic Factor Model Approach
    \\[0.4em]
    {\small Joint work with Ivan De Lorenzo Buratta (Prometeia)}
  \end{center}
\end{frame}

% ----------------------------------------------------------
% SLIDE 17: FC Motivation
% ----------------------------------------------------------
\begin{frame}{Essay III: Motivation}

\textbf{Why do we need financial cycle indicators?}
\begin{itemize}
  \item Cyclical systemic risk builds quietly — macroprudential policy must act \textbf{pre-emptively}
  \item The \textbf{Countercyclical Capital Buffer (CCyB)} requires quantitative benchmarks for activation
\end{itemize}

\vspace{0.5em}
\textbf{Current benchmark: credit-to-GDP gap (Basel III)}
\begin{itemize}
  \item HP-filter based — large \textbf{end-of-sample bias}; sensitive to output trends
  \item Uses only one variable: misses asset price and leverage dynamics
  \item Subject to \textbf{data vintage effects}: revised substantially post-crisis
\end{itemize}

\vspace{0.5em}
\textbf{This paper:} A Dynamic Factor Model that synthesizes \textbf{multiple financial indicators} into a single, statistically grounded financial cycle indicator.

\end{frame}

% ----------------------------------------------------------
% SLIDE 18: FC DFM Setup
% ----------------------------------------------------------
\begin{frame}{Essay III: DFM Framework}

\textbf{Country-specific DFM:}
\[
x_{it} = \lambda_i \, f_t + \varepsilon_{it}, \qquad i = 1,\ldots,p;\; t = 1,\ldots,T
\]

\begin{itemize}
  \item $x_{it}$: standardized financial indicator $i$ at time $t$ (credit growth, house prices, equity prices, leverage, current account, ...)
  \item $f_t$: latent \textbf{financial cycle factor} (scalar, I(1) or filtered to cyclical component)
  \item $\lambda_i$: factor loading — sign normalized so $f_t$ increases with systemic risk
  \item $\varepsilon_{it}$: idiosyncratic component, allowing weak cross-sectional dependence
\end{itemize}

\vspace{0.5em}
\textbf{Estimation:} Maximum likelihood / PCA with sign normalization.\\
\textbf{Data:} Quarterly, 1970--2022, for 10 European economies.

\end{frame}

% ----------------------------------------------------------
% SLIDE 19: FC Data and Variables
% ----------------------------------------------------------
\begin{frame}{Essay III: Data and Variable Selection}

\textbf{Financial indicators (country-specific panel):}

\vspace{0.4em}
\begin{tabular}{ll}
\toprule
\textbf{Category} & \textbf{Variables} \\
\midrule
Credit & Credit-to-GDP ratio, credit growth \\
Asset prices & House price growth, equity price growth \\
Leverage & Bank leverage, debt-service ratio \\
External & Current account balance \\
\bottomrule
\end{tabular}

\vspace{0.6em}
\textbf{Countries:} Austria, Belgium, Denmark, France, Germany, Ireland, Netherlands, Portugal, Spain, Sweden.

\vspace{0.5em}
\textbf{Key identification:} Factor loading $\lambda_i > 0$ for pro-cyclical indicators; sign normalization ensures $f_t$ tracks \textbf{systemic risk accumulation}.

\end{frame}

% ----------------------------------------------------------
% SLIDE 20: FC Country Results
% ----------------------------------------------------------
\begin{frame}{Essay III: Country-Specific Financial Cycle Indicators}

\textbf{Estimated financial cycle factors capture country-specific dynamics:}

\begin{itemize}
  \item \textbf{Ireland / Spain}: sharp pre-crisis buildup (2003--2007), pronounced decline post-2008
  \item \textbf{Germany}: relatively stable cycle; mild boom in 2000s
  \item \textbf{Portugal}: extended buildup through 1990s--2000s tied to euro accession credit expansion
  \item \textbf{Common element}: synchronized peak around 2007--2008 across all countries
\end{itemize}

\vspace{0.6em}
\textbf{Variable decomposition:} For each country and quarter, the factor loading $\hat{\lambda}_i$ reveals \textbf{which indicators} are driving the financial cycle signal at that point in time.

\end{frame}

% ----------------------------------------------------------
% SLIDE 21: FC Early Warning Properties
% ----------------------------------------------------------
\begin{frame}{Essay III: Early-Warning Properties}

\textbf{Crisis early-warning evaluation} (AUC-ROC, AUROC, 1970--2022 banking crises):

\vspace{0.4em}
\begin{tabular}{lc}
\toprule
\textbf{Indicator} & \textbf{AUROC (2-year horizon)} \\
\midrule
DFM financial cycle (this paper) & \textbf{0.82} \\
Credit-to-GDP gap (HP) & 0.68 \\
Credit growth & 0.70 \\
House price growth & 0.63 \\
\bottomrule
\end{tabular}

{\small AUROC values approximate from simulation exercises; higher = better early-warning capacity.}

\vspace{0.5em}
\textbf{Key result:} The DFM indicator \textbf{consistently outperforms} single-variable benchmarks. Signals the 2008 GFC with a 6--8 quarter lead across most countries.

\end{frame}

% ----------------------------------------------------------
% SLIDE 22: FC Neutral Risk Thresholds
% ----------------------------------------------------------
\begin{frame}{Essay III: Neutral Risk Environment}

\textbf{Beyond the indicator — defining a neutral risk benchmark:}

\vspace{0.3em}
A ``neutral'' CCyB environment corresponds to a financial cycle value in a \textbf{historically normal range}, not signaling elevated systemic risk.

\vspace{0.5em}
\textbf{Statistical approach:}
\begin{itemize}
  \item Estimate the empirical distribution of $\hat{f}_t$ in non-crisis periods
  \item Define the \textbf{neutral zone} as the central $[q_\alpha, q_{1-\alpha}]$ quantile interval
  \item \textbf{Buffer guidance}: zero CCyB inside neutral zone; graduated activation above $q_{1-\alpha}$
\end{itemize}

\vspace{0.5em}
\textbf{Advantage over subjective judgment:}
\begin{itemize}
  \item Data-driven, country-specific thresholds
  \item Transparent and reproducible
  \item Robust to the choice of $\alpha$ (sensitivity analysis available)
\end{itemize}

\end{frame}

% ----------------------------------------------------------
% SLIDE 23: FC Policy Implications
% ----------------------------------------------------------
\begin{frame}{Essay III: Policy Implications}

\textbf{Practical value for macroprudential authorities:}

\begin{itemize}
  \item \textbf{CCyB calibration}: DFM indicator provides a \textbf{quantitative benchmark} for buffer rate decisions, complementing the credit gap and macro-financial assessments
  \item \textbf{Early warning}: signals systemic risk buildup 6--8 quarters ahead — enough lead time for buffer accumulation
  \item \textbf{Variable decomposition}: identifies \textbf{which sector} (credit vs. real estate vs. leverage) is driving the cycle at any given moment, enabling targeted interventions
  \item \textbf{Cross-country comparability}: consistent methodology across 10 European economies enables \textbf{convergence monitoring} under the Single Supervisory Mechanism
\end{itemize}

\vspace{0.4em}
{\small Directly relevant for EBA/ECB macroprudential frameworks and ESRB risk assessment.}

\end{frame}

% ----------------------------------------------------------
% SLIDE 24: FC Takeaway
% ----------------------------------------------------------
\begin{frame}{Essay III: Key Takeaway}

\begin{itemize}
  \item The DFM approach provides a \textbf{statistically rigorous}, \textbf{multi-dimensional} measure of the financial cycle
  \item Outperforms the Basel III credit-to-GDP gap on early-warning criteria
  \item Country-specific indicators capture the \textbf{heterogeneity} of European financial cycles
  \item Neutral risk thresholds are \textbf{data-driven} and \textbf{transparent}, supporting operational CCyB guidance
  \item Framework is \textbf{extendable}: more countries, higher frequency, real-time updating
\end{itemize}

\end{frame}

% ==========================================================
% CONCLUSIONS
% ==========================================================

\begin{frame}{}
  \begin{center}
    \Large\textcolor{isegblue}{\textbf{Conclusions}}
  \end{center}
\end{frame}

% ----------------------------------------------------------
% SLIDE 25: Cross-Paper Contributions
% ----------------------------------------------------------
\begin{frame}{Contributions of the Thesis}

\begin{block}{Essay I — Smooth Transition Factor Model}
  First consistent estimator for factor models with \textbf{logistic smooth transitions} in loadings.
  Linearity test with correct size under $H_0$. Forecasting gains documented on FRED-QD.
\end{block}

\vspace{0.4em}

\begin{block}{Essay II — Nonstationary DFM with Threshold Cointegration}
  Tractable estimation framework for \textbf{threshold and band VECMs} in two-level factor models.
  Documents discontinuous adjustment in global yield curves consistent with friction theory.
\end{block}

\vspace{0.4em}

\begin{block}{Essay III — Financial Cycle Measurement}
  DFM-based financial cycle indicator with \textbf{superior early-warning capacity}.
  Country-specific, data-driven neutral thresholds for CCyB calibration across Europe.
\end{block}

\end{frame}

% ----------------------------------------------------------
% SLIDE 26: Methodological Advances
% ----------------------------------------------------------
\begin{frame}{Methodological Advances}

\textbf{Common thread:} All three essays \textbf{relax restrictive assumptions} in the standard linear DFM:

\vspace{0.5em}
\begin{tabular}{lll}
\toprule
\textbf{Essay} & \textbf{Assumption relaxed} & \textbf{Method} \\
\midrule
I   & Constant loadings       & Logistic smooth transition \\
II  & Linear adjustment       & Threshold / band VECM \\
III & Single indicator        & Multi-indicator latent factor \\
\bottomrule
\end{tabular}

\vspace{0.6em}
\textbf{Shared tools:}
\begin{itemize}
  \item PCA / SVD for initialization and factor extraction
  \item Grid search over nonlinear parameters
  \item Information criteria adapted for model complexity
  \item Monte Carlo validation before empirical application
\end{itemize}

\end{frame}

% ----------------------------------------------------------
% SLIDE 27: Future Research
% ----------------------------------------------------------
\begin{frame}{Future Research}

\begin{itemize}
  \item \textbf{Asymptotic distribution theory} for the smooth transition and threshold factor estimators (Essays I and II): confidence intervals and formal hypothesis tests
  \item \textbf{Impulse response functions} under nonlinear factor VECMs: regime-conditional transmission of global shocks
  \item \textbf{Multiple regimes}: extending Essays I and II to $K > 2$ regimes with endogenous regime selection
  \item \textbf{Real-time financial cycle monitoring}: extending Essay III to near-real-time updating with mixed-frequency data
  \item \textbf{Cross-country spillovers}: embedding country-specific DFM indicators into a global model of systemic risk contagion
\end{itemize}

\end{frame}

% ----------------------------------------------------------
% SLIDE 28: Thank You
% ----------------------------------------------------------
\begin{frame}{}

\begin{center}
  \vspace{1.5em}
  {\LARGE \textcolor{isegblue}{\textbf{Thank You}}}

  \vspace{1.5em}
  {\large Questions?}

  \vspace{2em}
  \rule{0.5\textwidth}{0.4pt}

  \vspace{1em}
  {\small Daniel Abreu\\
  ISEG -- Lisbon School of Economics and Management\\
  Universidade de Lisboa}
\end{center}

\end{frame}

% ==========================================================
% BACKUP SLIDES (not counted in 28)
% ==========================================================

\appendix

\begin{frame}{Backup: STFM — Factor Estimates}
\begin{center}
  \includegraphics[width=0.85\textwidth]{stfm_factors.pdf}
\end{center}
{\small Estimated factors from the STFM (black) vs.\ linear DFM (gray). Shaded areas: NBER recessions.}
\end{frame}

\begin{frame}{Backup: Financial Cycle — Variable Weights}
\textbf{Factor loadings $\hat{\lambda}_i$ by variable and country} (approximate):

\vspace{0.4em}
\begin{itemize}
  \item \textbf{Highest loadings}: house price growth, credit-to-GDP growth, bank leverage
  \item \textbf{Moderate loadings}: equity prices, debt-service ratio
  \item \textbf{Lowest loadings}: current account (stabilising role — often countercyclical)
\end{itemize}

\vspace{0.4em}
Loadings are \textbf{country-specific}: their relative importance varies with national financial structures.
\end{frame}

\begin{frame}{Backup: Essay II — Rotation Diagram}
\begin{center}
  \includegraphics[width=0.65\textwidth]{rotation_diagram.pdf}
\end{center}
{\small Two-level factor structure: global factor (top) connected to group-specific factors (middle) and observables (bottom).}
\end{frame}

\end{document}
